\documentclass[11pt]{article}
% Shared preamble for all daily exercises
% Update this file to change settings (padding, parindent, etc.) for all exercises

\usepackage[margin=0.75in]{geometry}
\usepackage{amsmath}
\usepackage{amssymb}

% Custom settings
\setlength{\parindent}{0pt}
\setlength{\parskip}{0.2cm}

\title{Daily Exercise}
\author{Dhruman Gupta}
\date{23rd April}

\begin{document}

% Common header for daily exercises
% This file can be included in other LaTeX documents

\maketitle

\noindent
\textbf{Course:} CS-3330, Theory of Computation

\vspace{0.2cm}

\noindent
\textbf{Question:}\noindent 
Show that if $P = NP$, then $EXP = NEXP$.

\textbf{Answer:}

Assume that

\[
P = NP.
\]

We will show that

\[
EXP = NEXP.
\]

We already know that

\[
EXP \subseteq NEXP,
\]

so it is enough to prove that $NEXP \subseteq EXP$.

Let $L \in NEXP$. Then there exists a nondeterministic TM $M$ and a polynomial $p$ such that $M$ decides $L$ in time $2^{p(n)}$ on inputs of length $n$.

Now define the padded language

\[
L' = \{\langle x, 1^{2^{p(|x|)}} \rangle \mid x \in L\}.
\]

We claim that $L' \in NP$. On input $\langle x, 1^{2^{p(|x|)}} \rangle$, an NP machine can guess an accepting computation of $M$ on $x$ and verify it step by step.

This is possible in polynomial time in the input length, because the input length is

\[
|x| + 2^{p(|x|)},
\]

and the guessed computation also has length at most $2^{p(|x|)}$. So the verification is polynomial in the size of the padded input. Hence,

\[
L' \in NP.
\]

Since $P = NP$, we get

\[
L' \in P.
\]

So there is a deterministic polynomial-time machine $A$ deciding $L'$. Let its running time be $m^c$ on inputs of length $m$.

Now we use $A$ to decide $L$. Given input $x$, construct the padded string

\[
\langle x, 1^{2^{p(|x|)}} \rangle
\]

and run $A$ on it.

The length of this input is

\[
m = |x| + 2^{p(|x|)}.
\]

Therefore the running time is at most

\[
m^c \leq (|x| + 2^{p(|x|)})^c \leq 2^{q(|x|)}
\]

for some polynomial $q$.

Hence this gives a deterministic exponential-time algorithm for $L$. Therefore,

\[
L \in EXP.
\]

Since $L$ was arbitrary, we have shown that

\[
NEXP \subseteq EXP.
\]

Combining this with $EXP \subseteq NEXP$, we get

\[
EXP = NEXP.
\]

Hence proved.

\end{document}
